
%! Author = wojtek
%! Date = 01.04.2026

% Preamble
\documentclass[11pt]{article}

% Packages
\usepackage{amsmath}
\usepackage{tabularx}
\usepackage{booktabs}
\usepackage{seqsplit}

% Document
\begin{document}

\begin{table}[h!]
    \centering
    \small
    \begin{tabularx}{\textwidth}{|l|X|X|}
        \hline
        File name & \seqsplit{aut\_branching} & \seqsplit{basic\_branching} \\
        \hline
		cubes3.grl & 0.12 & 0.06 \\
		cubes4.grl & 1.12 & 0.71 \\
		cubes5.grl & 23.61 & 17.29 \\
		trees36.grl & 0.24 & 0.16 \\
		wheeljoin14.grl & 9.97 & 7.76 \\
        \hline
    \end{tabularx}
    \caption{
        This run compares branching using generator-based automorphism counting v1\_0\_5 against basic branching v1\_0\_3 with naive automorphism counting.
        Note that the generator-based automorphism counting had a bug in the code during this run:
        it only considered a new branch to be an automorphism if it found a new generator.
        Thus, it explored almost the whole tree anyway, without any pruning.
    }
    \label{tab:1}
\end{table}

\end{document}
